% GENERATED FILE. Edit canonical lessons, then run npm run generate:books.
% canonical-source-hash: fnv1a64:5bb0b758bbf012ec
% canonical-lessons: RU-C164-pivo, RU-C164-vino, RU-C164-klimat, RU-C164-vozdukh, RU-C164-zemlya

\chapter{Beer, Wine, Climate, Air, Earth}
\label{ch:ru-beer-wine-climate-air-earth}
\hlchaptermodality{164}

\begin{chapteropening}
\textbf{By the end of this chapter:} \emph{I can say beer, wine, a climate, air and the earth.}

Complete the last lesson of chapter 164: I can say beer, wine, a climate, air and the earth.
\end{chapteropening}

\section[\texorpdfstring{\ru{пиво} (pívo)}{пиво (pívo)}]{\ru{пиво} (pívo) — beer}

\label{lesson:RU-C164-pivo}



\begin{quote}
Before the new one: say the Russian for cabbage soup, then the Russian for a pancake.
\end{quote}

\subsection*{You'll want to know: \ru{пиво}}
\textbf{\ru{пиво}} (\emph{pívo}) — \textquotedblleft{}beer\textquotedblright{}.

\begin{grammarlens}[title={What you've built}]
One more word for this chapter.
\end{grammarlens}

\subsection*{Your turn}
\begin{itemize}
\raggedright
  \item \emph{Say it:} \emph{pívo}
  \item \emph{Say it:} \emph{pívo}, once more
  \item \emph{Say it:} say \emph{pívo}
  \item \emph{Recall:} say the Russian for cabbage soup, then the Russian for a pancake, then say \emph{pívo} again
\end{itemize}

\begin{tcolorbox}[breakable,colback=teal!4,colframe=teal!35!black,title={Before you move on}]
What does \ru{пиво} mean? (\textquotedblleft{}Beer\textquotedblright{}.) Say it once more.
\end{tcolorbox}
\section[\texorpdfstring{\ru{вино} (vinó)}{вино (vinó)}]{\ru{вино} (vinó) — wine}

\label{lesson:RU-C164-vino}



\begin{quote}
Before the new one: say the Russian for a pancake, then the Russian for beer.
\end{quote}

\subsection*{You'll want to know: \ru{вино}}
\textbf{\ru{вино}} (\emph{vinó}) — \textquotedblleft{}wine\textquotedblright{}.

\begin{grammarlens}[title={What you've built}]
One more word for this chapter.
\end{grammarlens}

\subsection*{Your turn}
\begin{itemize}
\raggedright
  \item \emph{Say it:} \emph{vinó}
  \item \emph{Say it:} \emph{vinó}, once more
  \item \emph{Say it:} say \emph{vinó}
  \item \emph{Recall:} say the Russian for a pancake, then the Russian for beer, then say \emph{vinó} again
\end{itemize}

\begin{tcolorbox}[breakable,colback=teal!4,colframe=teal!35!black,title={Before you move on}]
What does \ru{вино} mean? (\textquotedblleft{}Wine\textquotedblright{}.) Say it once more.
\end{tcolorbox}
\section[\texorpdfstring{\ru{климат} (klímat)}{климат (klímat)}]{\ru{климат} (klímat) — a climate}

\label{lesson:RU-C164-klimat}



\begin{quote}
Before the new one: say the Russian for beer, then the Russian for wine.
\end{quote}

\subsection*{You'll want to know: \ru{климат}}
\textbf{\ru{климат}} (\emph{klímat}) — \textquotedblleft{}a climate\textquotedblright{}.

\begin{grammarlens}[title={What you've built}]
One more word for this chapter.
\end{grammarlens}

\subsection*{Your turn}
\begin{itemize}
\raggedright
  \item \emph{Say it:} \emph{klímat}
  \item \emph{Say it:} \emph{klímat}, once more
  \item \emph{Say it:} say \emph{klímat}
  \item \emph{Recall:} say the Russian for beer, then the Russian for wine, then say \emph{klímat} again
\end{itemize}

\begin{tcolorbox}[breakable,colback=teal!4,colframe=teal!35!black,title={Before you move on}]
What does \ru{климат} mean? (\textquotedblleft{}A climate\textquotedblright{}.) Say it once more.
\end{tcolorbox}
\section[\texorpdfstring{\ru{воздух} (vózdukh)}{воздух (vózdukh)}]{\ru{воздух} (vózdukh) — air}

\label{lesson:RU-C164-vozdukh}



\begin{quote}
Before the new one: say the Russian for wine, then the Russian for a climate.
\end{quote}

\subsection*{You'll want to know: \ru{воздух}}
\textbf{\ru{воздух}} (\emph{vózdukh}) — \textquotedblleft{}air\textquotedblright{}.

\begin{grammarlens}[title={What you've built}]
One more word for this chapter.
\end{grammarlens}

\subsection*{Your turn}
\begin{itemize}
\raggedright
  \item \emph{Say it:} \emph{vózdukh}
  \item \emph{Say it:} \emph{vózdukh}, once more
  \item \emph{Say it:} say \emph{vózdukh}
  \item \emph{Recall:} say the Russian for wine, then the Russian for a climate, then say \emph{vózdukh} again
\end{itemize}

\begin{tcolorbox}[breakable,colback=teal!4,colframe=teal!35!black,title={Before you move on}]
What does \ru{воздух} mean? (\textquotedblleft{}Air\textquotedblright{}.) Say it once more.
\end{tcolorbox}
\section[\texorpdfstring{\ru{земля} (zemlyá)}{земля (zemlyá)}]{\ru{земля} (zemlyá) — the earth; land}

\label{lesson:RU-C164-zemlya}



\begin{quote}
Before the new one: say the Russian for a climate, then the Russian for air.
\end{quote}

\subsection*{You'll want to know: \ru{земля}}
\textbf{\ru{земля}} (\emph{zemlyá}) — \textquotedblleft{}the earth; land\textquotedblright{}.

\begin{grammarlens}[title={What you've built}]
One more word for this chapter.
\end{grammarlens}

\subsection*{Your turn}
\begin{itemize}
\raggedright
  \item \emph{Say it:} \emph{zemlyá}
  \item \emph{Say it:} \emph{zemlyá}, once more
  \item \emph{Say it:} say \emph{zemlyá}
  \item \emph{Recall:} say the Russian for a climate, then the Russian for air, then say \emph{zemlyá} again
\end{itemize}

\begin{tcolorbox}[breakable,colback=teal!4,colframe=teal!35!black,title={Before you move on}]
What does \ru{земля} mean? (\textquotedblleft{}The earth; land\textquotedblright{}.) Say it once more.
\end{tcolorbox}
